% Glosario de notación del temario: símbolo, significado y dónde se usa.
% Empezado el 02/11/2026 al armonizar la notación de 3.A.10, 3.B.2, 3.B.4 y 3.B.23. Se amplía cuando se armoniza otro tema.
% Compilación (no se sube el PDF): latexmk -pdf -outdir=.build glosario.tex
\documentclass[10pt,a4paper]{article}
\usepackage[T1]{fontenc}
\usepackage[utf8]{inputenc}
\usepackage[spanish,es-noshorthands,es-nodecimaldot]{babel}
\usepackage{amsmath,amssymb}
\usepackage{array,longtable,booktabs}
\usepackage[margin=2cm]{geometry}
\usepackage[hidelinks]{hyperref}

\newcolumntype{S}{>{\raggedright\arraybackslash}p{2.6cm}}
\newcolumntype{M}{>{\raggedright\arraybackslash}p{9.2cm}}
\newcolumntype{D}{>{\raggedright\arraybackslash}p{3.7cm}}
\newcommand{\cabecera}{\toprule\textbf{Símbolo} & \textbf{Significado} & \textbf{Dónde}\\\midrule\endhead\bottomrule\endfoot}

\title{Glosario de notación}
\author{Temas armonizados: 3.A.10, 3.B.2, 3.B.4, 3.B.23 y 3.B.25\\Registrados, con diferencias pendientes: 3.B.13 y 3.B.24}
\date{2 de noviembre de 2026}

\begin{document}
\maketitle

\section*{Criterios}
\begin{enumerate}
    \item Dentro de cada tema, cada símbolo tiene un solo significado. Se exceptúan las letras que nombran los elementos de un ejemplo, como las acciones $A$, $B$ y $C$ o los proyectos $A$ y $B$, que se definen donde aparecen.
    \item Los conceptos que comparten varios temas se escriben con el mismo símbolo en todos ellos (apartado~1).
    \item Se conserva la notación original de un autor solo si cambiarla haría irreconocible un resultado conocido: es el caso de los parámetros de TVERSKY y KAHNEMAN en el 3.A.10.
    \item Las diferencias que aún quedan entre temas y las letras que se repiten dentro de un mismo tema se anotan al final como pendientes.
\end{enumerate}

\section{Notación común entre temas}
\begin{longtable}{SMD}
\cabecera
$r_f$ & Tipo de interés (rentabilidad) sin riesgo, en términos netos; $r_f^*$, el de la moneda extranjera. & 3.B.2, 3.B.4, 3.B.13 (interés simple), 3.B.23, 3.B.24, 3.B.25 (con capitalización continua)\\
$R_f$ & Tipo sin riesgo bruto, $1+r_f$. & 3.A.10, 3.B.23\\
$\theta$ & Coeficiente de aversión relativa al riesgo: parámetro de la utilidad CRRA y de la utilidad de media y varianza. & 3.A.10, 3.B.23; también 3.A.29, 3.A.38 y 3.A.44\\
$A(\cdot)$ & Coeficiente de aversión absoluta al riesgo, $-u''/u'$. & 3.A.10, 3.B.23\\
$R(\cdot)$ & Coeficiente de aversión relativa al riesgo, $-x\,u''(x)/u'(x)$; en el 3.B.23, $R(W)$. & 3.A.10, 3.B.23\\
$\delta$ & Factor de descuento temporal (impaciencia). & 3.A.10, 3.B.23\\
$u(\cdot)$ & Función de utilidad de VON NEUMANN y MORGENSTERN. & 3.A.10, 3.B.23\\
$\mu$, $\sigma$, $\sigma^2$, $\rho$ & Media, desviación típica (volatilidad), varianza y coeficiente de correlación. & Todos\\
$\beta$ & Beta: riesgo sistemático, covarianza con la cartera de mercado dividida por su varianza. & 3.B.2, 3.B.4, 3.B.23\\
$\lambda$ & Prima de riesgo: $\lambda_M$, del mercado; $\lambda_k$, del factor $k$; $\lambda=\beta\lambda_M$, la de un activo en el CAPM. & 3.B.2, 3.B.23\\
$\alpha$ & Alfa: la parte de la rentabilidad esperada que el modelo no explica (constante de la regresión). & 3.B.4, 3.B.23\\
$\varepsilon$ & Componente específico (idiosincrático) de una rentabilidad, de media nula. & 3.B.2, 3.B.4, 3.B.23\\
$F_k$ & Sorpresa del factor $k$, de media nula. & 3.B.2, 3.B.23\\
$E$, $D$, $V$ & Valores de mercado de los recursos propios, de la deuda y de la empresa, $V=E+D$. & 3.B.2, 3.B.4\\
$M$ & Cartera de mercado. & 3.B.2, 3.B.23\\
$B(0,t)$ & Precio hoy de un bono cupón cero sin riesgo que paga una unidad en $t$: el factor de descuento de ese plazo. El 3.B.2 escribe $B(t,s)$, su precio en $t$ con pago en $s$. & 3.B.2, 3.B.24, 3.B.25\\
$S$, $F$ & Tipo de cambio al contado y a plazo, en unidades de moneda nacional por unidad de la extranjera (cotización directa): $S$ sube cuando la moneda nacional se deprecia. & 3.B.13, 3.B.25\\
$x$ & Base de divisas: tipo directo del dólar menos el tipo sintético que se obtiene con un \textit{FX swap}, $x=r_f^*-r^{\$}_{\mathrm{sint}}$; negativa cuando los dólares sintéticos son más caros. & 3.B.13, 3.B.25 (con capitalización continua)\\
$K$, $c$, $p$ & Precio de ejercicio y primas de la \textit{call} y de la \textit{put} europeas. & 3.B.13, 3.B.25\\
\end{longtable}

\section{Pendientes de armonizar}
\subsection*{Diferencias entre temas}
\begin{longtable}{SMD}
\toprule\textbf{Concepto} & \textbf{Símbolos} & \textbf{Temas}\\\midrule\endhead\bottomrule\endfoot
Estado de la naturaleza & $\omega$ en el 3.A.10 y el 3.B.25; $s$ en el 3.B.23, donde $\omega$ son los pesos de una cartera. & 3.A.10, 3.B.23, 3.B.25\\
Pesos de una cartera & $w_i$ en el 3.B.2; $\omega_j$ en el 3.B.23. & 3.B.2, 3.B.23\\
Riqueza & $w$ (y $x$, como resultado) en el 3.A.10; $W$ en el 3.B.23. & 3.A.10, 3.B.23\\
Esperanza & $E[\cdot]$ en el 3.A.10, el 3.B.24 y el 3.B.25 (en este, $E^{\mathbb{Q}}[\cdot]$ con la medida neutral al riesgo); $\mathbb E[\cdot]$ en el 3.B.2, el 3.B.13 ($\mathbb{E}_t[\cdot]$, condicionada a la información de $t$) y el 3.B.23. & 3.A.10, 3.B.2, 3.B.13, 3.B.23, 3.B.24, 3.B.25\\
Factor de descuento estocástico & $m_{t+1}$ en el 3.B.23 y el 3.B.24 (en este, $m$ es también el número de cupones al año); en el 3.B.25, el cociente $\psi_\omega/\pi_\omega$ entre el precio de estado y la probabilidad real, sin letra propia, porque $m$ es allí el diferencial de un préstamo sobre el tipo de referencia. & 3.B.23, 3.B.25\\
Letra $\pi$ & Prima de riesgo de PRATT en el 3.A.10; inflación esperada en el 3.B.2; probabilidad de un estado en el 3.B.23; inflación, $\pi$, e inflación esperada, $\pi^e$, en el 3.B.24; prima de riesgo de cambio, $\pi_t$, en el 3.B.13. & 3.A.10, 3.B.2, 3.B.13, 3.B.23, 3.B.24\\
Rentabilidad & $r$ minúscula en el 3.B.2 y el 3.B.23 (en el 3.B.23, $R$ es la bruta); $R_{it}$ en los estudios de acontecimientos del 3.B.4. & 3.B.2, 3.B.4, 3.B.23\\
Rentabilidad esperada & $\mu_j$ en el 3.B.23; $\mathbb E[r_i]$ en el 3.B.2. Índice del título: $j$ en el 3.B.23, $i$ en el 3.B.2. & 3.B.2, 3.B.23\\
Rentabilidad exigida o coste del capital & $r$ en el 3.B.2; $k_B$ en el 3.B.4; $k_j$ en el 3.B.23 (descuento de dividendos). & 3.B.2, 3.B.4, 3.B.23\\
Tipo impositivo & $T$ en el 3.B.2; $\tau$ en el 3.B.4, el 3.B.23 y el 3.B.24. & 3.B.2, 3.B.4, 3.B.23, 3.B.24\\
Equivalente cierto & $EC$ y $c(F,u)$ en el 3.A.10, dos notaciones en el mismo tema; en el 3.B.23, $U$ es el equivalente cierto de la utilidad de media y varianza. & 3.A.10, 3.B.23\\
Letra $\delta$ & Factor de descuento temporal en el 3.A.10 y el 3.B.23; rendimiento continuo del subyacente (dividendos, intereses de una divisa) en el 3.B.4 y el 3.B.25. Aceptado por el usuario: son conceptos distintos. & 3.A.10, 3.B.4, 3.B.23, 3.B.25\\
Letra $\lambda$ & Prima de riesgo en el 3.B.2 y el 3.B.23; en el 3.B.25, precio de mercado del riesgo, $(\mu-r_f)/\sigma$, la prima por unidad de volatilidad (las intensidades de impago y de saltos son allí $\gamma$ y $\eta$); intensidad de impago en el 3.B.24; en el 3.B.13, cuánto mueve el precio cada unidad de flujo de órdenes (KYLE, que lo llama así). & 3.B.2, 3.B.13, 3.B.23, 3.B.24, 3.B.25\\
Letra $\gamma$ & Parámetro de la función de ponderación de probabilidades en el 3.A.10; persistencia del beneficio residual (OHLSON) en el 3.B.23; intensidad de impago de un CDS en el 3.B.25. & 3.A.10, 3.B.23, 3.B.25\\
Beneficio y valor contable & $B_1$ es el beneficio por acción en el 3.B.2 ($B_1/P_0$, inversa del PER); en el 3.B.23, $Y_t$ es el beneficio y $B_t$ el valor contable de los recursos propios. & 3.B.2, 3.B.23\\
Siglas IR & Índice de rentabilidad de una inversión en el 3.B.2; ratio de información en el 3.B.23 (ley fundamental de GRINOLD). & 3.B.2, 3.B.23\\
Letra $F$ & Sorpresa del factor $k$, $F_k$, en el 3.B.2 y el 3.B.23; precio a plazo, $F_t$, en el 3.B.25, y tipo de cambio a plazo, $F$, en el 3.B.13. & 3.B.2, 3.B.13, 3.B.23, 3.B.25\\
Impago & Pérdida en caso de impago: $L$ en el 3.B.2; $\ell$ en el 3.B.24, donde $L$ son los pasivos de un balance; el 3.B.25 usa la fracción del nominal que se recupera, $R$, y la pérdida es $1-R$. En el modelo de intensidad del 3.B.24 (DUFFIE y SINGLETON, 1999), $\ell$ es una fracción del valor de mercado y no del nominal, de modo que $s=\lambda\ell$ y $s=\gamma(1-R)$ solo coinciden de forma aproximada. Intensidad de impago: $\lambda$ en el 3.B.24 y $\gamma$ en el 3.B.25. & 3.B.2, 3.B.24, 3.B.25\\
Tipo a plazo & $r_{f,t,s}$ en el 3.B.2; $f_{t,s}$ en el 3.B.24; $L^F_i$, el del periodo $i$ de un FRA o de un swap, en el 3.B.25. & 3.B.2, 3.B.24, 3.B.25\\
Letra $y$ & Rendimiento al vencimiento (TIR) en el 3.B.2 y el 3.B.24; rendimiento de conveniencia en el 3.B.25. & 3.B.2, 3.B.24, 3.B.25\\
Letra $\beta$ & Beta del riesgo sistemático en el 3.B.2, el 3.B.4 y el 3.B.23; parámetros de NELSON y SIEGEL, $\beta_0$, $\beta_1$ y $\beta_2$, en el 3.B.24, que son los del autor; coeficiente de la regresión de FAMA en el 3.B.13, también el del autor, donde la exposición cambiaria de ADLER y DUMAS es $\beta_S$. & 3.B.2, 3.B.4, 3.B.13, 3.B.23, 3.B.24\\
Letra $D$ & Deuda en el 3.B.2 y el 3.B.4 (en el 3.B.2, también duración); duración de MACAULAY en el 3.B.24, donde la deuda de un balance son los pasivos, $L$; intermediario del ejemplo conductor en el 3.B.13. & 3.B.2, 3.B.4, 3.B.13, 3.B.24\\
Distribución normal & $N(\cdot)$ en el 3.B.25; $\Phi(\cdot)$ en el 3.B.13, donde $N$ es el nominal en dólares (también es nominal en el 3.B.24). & 3.B.13, 3.B.24, 3.B.25\\
Capitalización de $r_f$ y $r_f^*$ & Interés simple del mercado monetario en el 3.B.13, con la paridad cubierta $F=S\,(1+r_f\tau)/(1+r_f^*\tau)$; continua en el 3.B.25, $F_0=S_0\,e^{(r_f-r_f^*)T}$. El 3.B.13 usa la forma continua solo con GARMAN y KOHLHAGEN. & 3.B.13, 3.B.25\\
Tipo de cambio y tipos de interés en el 3.B.14 & $E$ para el contado e $i$, $i^*$ para los tipos de interés, frente a $S$, $r_f$ y $r_f^*$ en el 3.B.13 y el 3.B.25; $E$ choca allí con la esperanza. Se armoniza cuando se trabaje el 3.B.14. & 3.B.13, 3.B.14, 3.B.25\\
\end{longtable}

\subsection*{Letras que se repiten dentro de un tema}
\begin{longtable}{SMD}
\toprule\textbf{Tema} & \textbf{Letras y sus sentidos} & \textbf{Gravedad}\\\midrule\endhead\bottomrule\endfoot
3.A.10 & $L$: lotería y pérdida del seguro. $E$: suceso (utilidad de CHOQUET) y esperanza. $a$: acto y coeficiente de la CARA. $\sigma$: desviación típica y elasticidad de sustitución intertemporal. $w$: riqueza y función de ponderación de TVERSKY y KAHNEMAN. $c$: consumo, equivalente cierto $c(F,u)$ y precio de la información $c^*$. $\alpha$: peso de la mezcla y curvatura de TVERSKY y KAHNEMAN (conservada). & Media: apartados distintos\\
3.B.2 & $T$: vida del proyecto y tipo impositivo. $D$: deuda, duración y dividendo $D_{i,1}$. $d_1$: argumento de BLACK y SCHOLES y dividendo del próximo año. $P$: periodo de recuperación, precio y proyecto. $p$: probabilidad y precio en GORDON y SHAPIRO. $C$: coste de restauración, valor de la opción y empresa comparable. & Media\\
3.B.4 & $P$: precio pagado por la empresa y precio de ejercicio de la opción real. $q$: probabilidad neutral al riesgo; $Q_A$ y $Q_B$, valoraciones de las empresas. & Baja\\
3.B.13 & $x$: base de divisas (bloques 1 y 3, anexos I.5 y II.6) y orden del inversor informado de KYLE (II.3). $m$: precio medio del par (I.2) y valor que el intermediario espera de la divisa en el modelo de inventario (II.2). $c$: prima de la \textit{call} (I.6), cupón en euros de un swap de divisas, con $c^*$ el de dólares (figura de I.5), y $c_{j,t}$, unidades de la moneda $j$ por unidad de la nacional en el índice efectivo (I.2). $p$: prima de la \textit{put} (I.6) y precio $p$, $p_0$, de KYLE (II.3). $V$: valor de la divisa en los modelos de información, $V_H$ y $V_L$ (II.3), y valor de la empresa, $V_t$, en ADLER y DUMAS (I.7). & Baja\\
3.B.23 & $a$, $b$, $c$, $d$: constantes de la frontera y, según el apartado, peso del título en la tangencia ($a$), carteras de la APT ($a$, $b$), cargas ($b_j$), consumo ($c_t$) y valor del derecho de suscripción ($d$). $P$: precio, cartera factorial $P_k$ y activo $P$ mal valorado. $s$: carga en SMB y estado de la naturaleza. $S$: número de estados; $S_p$, ratio de SHARPE. $x$: peso en la APT con un factor y pago en el factor de descuento. $k$: rentabilidad exigida ($k_j$) e índice de los factores y de los retardos ($\lambda_k$, $\rho_k$). $\mu$, $\sigma$: media y volatilidad de las rentabilidades y, en la desigualdad de CHEBYSHEV (III.3.3), media y desviación típica del precio. $b$: además, tasa de retención (III.4). $g$: crecimiento del consumo (II.7.4) y de los dividendos (I.1.2, III.4). $x_t$: logaritmo de la rentabilidad por dividendo (III.4.6). $A$: aversión absoluta, $A(\cdot)$, y cartera activa de TREYNOR y BLACK (III.4.8). $H$: horizonte (I.1.2, III.4); la historia de precios es $\mathcal H_t$ (III.3). $V$: valor de los activos ($V_T$, $V_M$) y valor que estima el analista ($V_t$, III.4). $\phi$ y $\varphi$: factor de CAMPBELL y SHILLER (III.4.6) y multiplicadores de LAGRANGE ($\varphi_1$, $\varphi_2$). $D$: deuda, en el pago del accionista, $\text{máx}(V-D,0)$ (I.1.1 y bloque 1 del cuerpo), y dividendo, $D_{j,t}$. $A$ y $B$: además, las dos acciones del ejemplo del cuerpo. & Baja: definidas donde aparecen\\
3.B.24 & $s$: diferencial de crédito (bloque 1, I.4) y segunda fecha del tipo a plazo $f_{t,s}$ (bloque 3, I.9). $\tau$: fracción transcurrida del periodo de cupón (I.2), tipo impositivo (I.5) y parámetro de NELSON y SIEGEL (bloque 3, I.10, I.14). $m$: cupones al año (I.1-I.6) y factor de descuento estocástico, $m_{t+1}$ (I.14). $k$: fecha de un cambio del tipo (bloque 1), índice de los pagos en la fórmula del Tesoro (I.2) y año futuro del tipo $i_k$ (bloque 3, I.12). $E$: capital de un balance, $E=A-L$ (bloque 1, I.8), y esperanza, $E[\cdot]$ (bloque 3, I.12-I.14). $i$: tipo corto (bloque 3, I.12-I.14) y tipo de descuento de las letras, $i_d$ (I.3). & Baja: primera pasada; definidas donde aparecen\\
3.B.25 & Resueltas el 04/11/2026 (segunda pasada): cada variable del tema tiene un solo símbolo. Quedan, por ser letras de un modelo concreto dentro de su apartado, $u$ y $d$ (factores del árbol) junto a la utilidad $u_A$ y a $d_1$, $d_2$ de BLACK-SCHOLES; $I$ (rentas) y la inversión $I$ de FROOT, SCHARFSTEIN y STEIN, con su coste $C(\cdot)$ junto a la \textit{call} americana $C$; $\theta$ y $\rho$ de HESTON junto a la aversión al riesgo y a la rho; $\omega$ del GARCH junto al estado $\omega$; $k$ de LELAND y $k_A$, $k_B$ de la ventaja comparativa; $\varepsilon$ del GARCH junto a la normal estándar; $e_1$, $e_2$ junto a la dotación $e$, y $x$, genérica, junto a la base de divisas. & Baja\\
\end{longtable}

\section{3.B.23: Análisis de los instrumentos financieros de renta variable}
\subsection*{Anexo I: renta variable y elección de cartera}
\begin{longtable}{SMD}
\cabecera
$r_j$, $r_{j,t}$ & Rentabilidad de la acción $j$ (en el periodo $t$). & I.2\\
$r^c$ & Rentabilidad logarítmica o continua, $\ln(1+r)$. & I.2.1\\
$P_{j,t}$, $D_{j,t}$ & Precio y dividendo por acción en $t$. & I.1, I.2\\
$k_j$ & Rentabilidad que exigen los inversores a la acción $j$ en el descuento de dividendos. & I.1.2\\
$d$ & Valor teórico del derecho de suscripción. & I.1.3\\
$\bar r_a$, $\bar r_g$ & Medias aritmética y geométrica de las rentabilidades. & I.2.2\\
$\hat U_H$ & Estimador de BLUME de la rentabilidad media a $H$ periodos. & I.2.2\\
$W_t$ & Riqueza en $t$. & I.2, I.7, I.9\\
$\mu_j=\mathbb E[r_j]$ & Rentabilidad esperada de la acción $j$. & Todo el tema\\
$\sigma_j$, $\sigma_{jk}$, $\rho_{jk}$ & Volatilidad; covarianza y correlación entre $j$ y $k$. & I.3\\
$\omega_j$, $\boldsymbol\omega$ & Peso de la acción $j$ en una cartera; vector de pesos. & I.3-I.9, II.5\\
$\boldsymbol\mu$, $\boldsymbol\Sigma$, $\mathbf 1$ & Vector de rentabilidades esperadas, matriz de covarianzas y vector de unos. & I.3.3-I.9\\
$r_p$, $\mu_p$, $\sigma_p$ & Rentabilidad, rentabilidad esperada y volatilidad de la cartera $p$. & I.3-II\\
$\sigma_{jp}$, $\beta_{jp}$ & Covarianza de la acción con la cartera $p$ y beta respecto de ella. & I.3.3\\
$\bar\sigma^2$, $\bar\sigma_{jk}$, $\bar\rho$ & Varianza, covarianza y correlación medias. & I.4.1\\
$F$, $b_j$ & Factor común y sensibilidad de la acción a él (modelo de un factor); en el cuerpo, $\beta_j$. & I.4.2\\
$\varepsilon_j$, $\sigma_{\varepsilon_j}^2$ & Componente específico y su varianza. & I.4.2, I.8, II.5\\
$\varphi_1$, $\varphi_2$ & Multiplicadores de LAGRANGE del problema de MARKOWITZ. & I.5.2, II.1.3\\
$a$, $b$, $c$, $d$ & Constantes de la frontera de MERTON (1972): $a=\mathbf 1'\boldsymbol\Sigma^{-1}\boldsymbol\mu$, $b=\boldsymbol\mu'\boldsymbol\Sigma^{-1}\boldsymbol\mu$, $c=\mathbf 1'\boldsymbol\Sigma^{-1}\mathbf 1$, $d=bc-a^2$. & I.5.2, I.7.2, II.3.1\\
$MV$ & Cartera de mínima varianza global. & I.5\\
$C_1$, $C_2$ & Dos carteras de la frontera que hacen de fondos en la propiedad de los dos fondos (BLACK, 1972); $\omega$, peso del primero. & Cuerpo, 2.4.1\\
$T$ & Cartera tangente, la de mayor ratio de SHARPE. & I.6-II\\
$S_p$ & Ratio de SHARPE de la cartera $p$, $(\mu_p-r_f)/\sigma_p$. & I.6.2, II.2\\
$\mathbf z$ & $\boldsymbol\Sigma^{-1}(\boldsymbol\mu-r_f\mathbf 1)$: pesos sin normalizar de la cartera tangente. & I.6.3\\
$y$, $y^*$ & Proporción de la riqueza en la cartera con riesgo y proporción óptima. & I.6-I.7\\
$U$ & Utilidad de media y varianza, $\mu_p-(\theta/2)\,\sigma_p^2$: equivalente cierto de la cartera. & I.7\\
$\alpha_j$, $\beta_j$ & Constante y pendiente del modelo de índice único. & I.8\\
$N$, $K$ & Número de acciones y de factores. & Todo el tema\\
\end{longtable}

\subsection*{Anexo II: valoración}
\begin{longtable}{SMD}
\cabecera
$r_M$, $\mu_M$, $\sigma_M$ & Rentabilidad, rentabilidad esperada y volatilidad de la cartera de mercado. & II\\
$\sigma_{jM}$, $\beta_j$ & Covarianza con el mercado y beta, $\sigma_{jM}/\sigma_M^2$. & II\\
$\omega_{jM}$ & Peso de la acción $j$ en la cartera de mercado; $\sum_j\omega_{jM}\beta_j=1$. & Cuerpo, 3.1.1\\
$A'$, $B'$ & Combinaciones del mercado con el activo sin riesgo con la misma beta que $A$ y $B$ (CML frente a SML). & Cuerpo, 3.1.1\\
$\lambda_M$ & Prima de riesgo del mercado, $\mu_M-r_f$. & II\\
$\bar\theta$ & Aversión relativa media del mercado: media armónica de las de los inversores, ponderada por su riqueza. & II.1.2\\
$\Delta\omega$ & Pequeño aumento del peso de una cartera o de una acción. & II.1.3\\
$X$, $P_0$ & Pago incierto y precio de un activo (equivalente cierto). & II.1.5\\
$\alpha_j$ & Alfa: rentabilidad esperada por encima de la que exige el modelo. & II.1.5, II.2\\
$\mathit{TR}_p$ & Ratio de TREYNOR, $(\bar r_p-r_f)/\beta_p$. & II.2.1\\
$M^2$ & Medida de MODIGLIANI y MODIGLIANI. & II.2.1\\
$Z$, $\mu_Z$ & Cartera de beta nula y su rentabilidad esperada. & II.3.1\\
$\tau$ & Tipo efectivo de BRENNAN: diferencia entre la tributación de dividendos y plusvalías. & II.3.2\\
$D_j/P_j$ & Rentabilidad por dividendo. & II.3.2\\
$\beta_{jk}$ & Sensibilidad de la acción $j$ al factor $k$. & II.3.3, II.5, II.6\\
$\beta_{jc}$, $\lambda_c$, $\Delta c$ & Beta de consumo, su prima y crecimiento del consumo. & II.3.3, II.7\\
$J$ & Estadístico del contraste de GIBBONS, ROSS y SHANKEN. & II.4.1\\
$\lambda_{k,t}$ & Coeficientes de la regresión mensual de FAMA y MacBETH: primas estimadas cada mes. & II.4.1\\
$\sigma_{\varepsilon_p}$, $\eta_{p,t}$ & Riesgo específico de la cartera y residuo de esa regresión. & II.4.1\\
$\lambda_0$, $\lambda_k$ & Constante y prima del factor $k$ en la APT. & II.5\\
$P_k$, $r_{P_k}$ & Cartera factorial pura del factor $k$ y su rentabilidad. & II.5.3\\
$x$ & Peso de la cartera $a$ en la combinación sin riesgo con un factor. & II.5.3\\
$r_R$, $\mu_R$ & Rentabilidad de la réplica y su esperanza. & II.5.3\\
$\alpha_j$ (APT) & Error de valoración respecto de la relación de la APT. & II.5.4\\
SMB, HML, RMW, CMA, Mom & Factores de tamaño, valor, rentabilidad, inversión e impulso. & II.6\\
$b_j$, $s_j$, $h_j$ & Cargas en el mercado, SMB y HML del modelo de tres factores; las de una cartera, $b_p$, $s_p$, $h_p$. & II.6.1 y cuerpo, 3.3.1\\
$\lambda_{\text{SMB}}$, $\lambda_{\text{HML}}$, $\sigma_{\text{SMB}}$, $\sigma_{\text{HML}}$ & Primas (rentabilidades esperadas) y volatilidades de SMB y HML. & II.6.1 y cuerpo, 3.3.1\\
$M_t$, $B_t$, $Y$, $dB$ & Valor de mercado, valor contable, beneficio e inversión en la ecuación de FAMA y FRENCH (2015); en el cuerpo, por acción: $P_{j,0}$, $B_{j,t}$, $Y_{j,t}$, $\Delta B_{j,t}$, con la tasa $k_j$. & II.6.2 y cuerpo, 3.3.2\\
$m$ & Factor de descuento estocástico. & II.7 y cuerpo, 3.3.4\\
$c_t$, $\xi$ & Consumo y cantidad comprada del activo. & II.7.1\\
$p$, $x$ & Precio y pago de un activo. & II.7 y cuerpo, 3.3.4\\
$R$, $R^e$ & Rentabilidad bruta y rentabilidad en exceso. & II.7\\
$\beta_{j,m}$, $\lambda_m$ & Beta respecto del factor de descuento y precio del riesgo. & II.7.1\\
$s$, $S$ & Estado de la naturaleza y número de estados. & II.7.2\\
$\pi_s$, $q_s$, $m_s$ & Probabilidad, precio del estado y factor de descuento en el estado $s$. & II.7.2\\
$\pi_s^*$, $\mathbb E^*$ & Probabilidad neutral al riesgo y esperanza con ella. & II.7.2\\
$\mathbf f$, $\mathbf b$, $\boldsymbol\Sigma_f$, $\boldsymbol\lambda$ & Factores y coeficientes de un factor de descuento lineal, covarianzas de los factores y primas. & II.7.3\\
$g$, $\sigma_c$ & Media y volatilidad del crecimiento del consumo. & II.7.4\\
\end{longtable}

\subsection*{Anexo III: eficiencia y análisis de inversiones}
\begin{longtable}{SMD}
\cabecera
$\Phi_t$ & Conjunto de información disponible en $t$. & III.1.1 y cuerpo, 4.1\\
$k_{j,t+1}$, $k$ & Rentabilidad exigida a la acción $j$ entre $t$ y $t+1$, la $k_j$ del anexo I; $k$, si es constante. En el cuerpo (4.1), $k_{t+1}$, sin el subíndice de la acción. & III.1.1, III.3.1 y cuerpo, 4.1\\
$\nu$, $e_t$ & Deriva y perturbación de un paseo aleatorio en precios. & III.1.2\\
$\mathbb E_t[\cdot]$ & Esperanza con la información disponible en $t$. & III.1.2\\
$\rho_k$ & Autocorrelación de las rentabilidades a $k$ periodos. & III.1.3\\
$VR(q)$ & Razón de varianzas a $q$ periodos. & III.1.3\\
$\hat\varepsilon_{j,t}$ & Rentabilidad anormal de la acción $j$ en $t$. & III.1.3\\
$\kappa$, $\psi(\kappa)$, $\chi$ & Proporción de inversores informados, ventaja de estar informado y coste de informarse (GROSSMAN y STIGLITZ). & III.2.1\\
$\mathcal H_t$ & Historia de precios y volúmenes hasta $t$, la información de la forma débil. & III.3.1 y cuerpo, 4.2\\
$\omega_t=\omega(\mathcal H_t)$ & Peso que una regla técnica da al activo según esa historia. & III.3.1 y cuerpo, 4.2\\
$n$ & Número de sesiones de una media móvil o de un indicador. & III.3.3\\
$\text{SMA}_t(n)$, $\text{EMA}_t(n)$, $\zeta$ & Medias móviles simple y exponencial de $n$ sesiones y factor de suavizado de la exponencial, habitualmente $2/(n+1)$. & III.3.3\\
$\text{MACD}_t$ & Diferencia entre las medias exponenciales de $12$ y $26$ sesiones. & III.3.3\\
$G_t$, $L_t$, $\text{RSI}_t$ & Medias suavizadas de las subidas y de las bajadas e índice de fuerza relativa. & III.3.3\\
$\hat\sigma_t(n)$ & Desviación típica de los $n$ últimos cierres (bandas de BOLLINGER). & III.3.3\\
$\ell$ & Número de desviaciones típicas en la desigualdad de CHEBYSHEV. & III.3.3\\
$V_t$, $P_t$ & Valor que estima el analista y precio de la acción. & III.4.1 y cuerpo, 4.3.1\\
$H$ & Horizonte: periodos hasta que el precio alcanza el valor estimado (III.4.1) y años de previsión explícita (III.4.3), como en el anexo I. & III.4.1, III.4.3\\
$Y_t$, $B_t$, $D_t$ & Beneficio del ejercicio $t$, valor contable de los recursos propios al final de $t$ y dividendo, por acción. & III.4\\
$\text{ROE}_t$ & Rentabilidad de los recursos propios, $Y_t/B_{t-1}$. & III.4.2-III.4.5\\
$b$, $g$, $g_L$ & Tasa de retención, crecimiento ($g=b\cdot\text{ROE}$, el sostenible) y crecimiento a largo plazo. & III.4.3, III.4.4\\
VAOC & Valor actual de las oportunidades de crecimiento, $P_0-Y_1/k$. & III.4.3, III.4.4 y cuerpo, 4.3.2\\
PER, P/VC & Ratios entre precio y beneficio y entre precio y valor contable. & III.4.4\\
$Y^a_t$, $\gamma$ & Beneficio residual, $Y_t-k\,B_{t-1}$, y su persistencia en el modelo de OHLSON (1995). & III.4.5\\
$x_t$, $\phi$ & Logaritmo de la rentabilidad por dividendo y factor de la aproximación de CAMPBELL y SHILLER (1988a), $1/(1+e^{\bar x})$. & III.4.6\\
CAPE & Precio real entre la media del beneficio real de los diez años anteriores. & III.4.6\\
$A$, $\alpha_A$, $\beta_A$, $\sigma_{\varepsilon_A}$ & Cartera activa de TREYNOR y BLACK (1973), su alfa, su beta y su riesgo específico. & III.4.8 y cuerpo, 4.3.4 ($A$)\\
$z_M$, $z_A$, $\omega_A^0$, $\omega_A^*$ & Pesos sin normalizar del mercado y del alfa puro, peso de la cartera activa sin corregir por su beta y peso óptimo. & III.4.8\\
$T^*$, $S_{T^*}$ & Combinación óptima del mercado y la cartera activa y su ratio de SHARPE. & III.4.8 y cuerpo, 4.3.4 ($T^*$)\\
IR, IC, BR & Ratio de información, coeficiente de información y amplitud (ley fundamental de GRINOLD). & III.4.8\\
\end{longtable}

\section{3.B.2: La empresa y las decisiones de inversión}
\begin{longtable}{SMD}
\cabecera
$r$ & Tasa de descuento o coste del capital de una inversión, $r=r_f+\lambda$. & Todo el tema\\
$r_{f,t}$, $r_{f,t,s}$ & Tipo sin riesgo en $t$ y para el plazo entre $t$ y $s$. & Bloque 1\\
$B(t,s)$ & Precio en $t$ de un bono cupón cero que paga una unidad en $s$. & Bloque 1\\
$i$, $\pi$ & Tipo nominal e inflación esperada (FISHER). & Bloque 1\\
$CF_t$, $I_0$, $T$ & Flujo de caja neto del periodo $t$, inversión inicial y vida del proyecto. & Bloque 1\\
$VAN$, $VA$, $TIR$, $IR$ & Valor actual neto, valor actual, tasa interna de retorno e índice de rentabilidad. & Bloque 1\\
$I_A$, $CF_A$, $I_B$, $CF_B$ & Inversión y flujo de los proyectos $A$ y $B$ del ejemplo. & Bloque 1\\
$r_F$ & Tasa de FISHER, la que iguala los VAN de dos proyectos. & Bloque 1\\
$K$ & Presupuesto disponible (racionamiento de capital). & Bloque 1\\
$P$, $P_d$ & Periodo de recuperación simple y descontado. & Bloque 1\\
$C$ & Coste de restaurar el terreno (ejemplo de varias TIR). & Bloque 1\\
$I^*$ & Inversión óptima en la separación de FISHER. & Bloque 1\\
$p$, $p^*$ & Probabilidad del escenario favorable y probabilidad crítica (punto muerto). & Bloque 1\\
$CF^+$, $CF^-$ & Flujos en el escenario favorable y en el desfavorable. & Bloque 1\\
$C$, $S^*$, $V(I_0)$ & Valor de la opción real de esperar, valor de los flujos sin los que se pierden por esperar y valor actual de la inversión pospuesta. & Bloque 1\\
$\sigma$, $t$, $N(\cdot)$, $d_1$, $d_2$ & Volatilidad del valor del proyecto, plazo hasta la decisión, distribución normal estándar y argumentos de BLACK y SCHOLES. & Bloque 1\\
$\lambda$, $\lambda_M$, $\lambda_E$, $\lambda_D$ & Prima de riesgo: de un activo, del mercado, de las acciones y de la deuda. & Bloque 2\\
$\beta_E$, $\beta_A$, $\beta_D$, $\beta_P$ & Betas de las acciones, de los activos, de la deuda y del proyecto. & Bloque 2\\
$r_E$, $r_D$, $r_A$, $r_P$ & Costes de los recursos propios, de la deuda, de los activos y del proyecto. & Bloque 2\\
$r_{\text{WACC}}$ & Coste medio ponderado del capital. & Bloque 2\\
$T$ & Tipo del impuesto de sociedades (en el WACC). & Bloque 2\\
$\,^C$, $\,^P$ & Superíndices de la empresa comparable y del proyecto. & Bloque 2\\
$y$, $p$, $L$ & Rentabilidad al vencimiento, probabilidad de impago y pérdida en caso de impago: $r_D\approx y-p\,L$. & Bloque 2\\
$d_1$, $p$, $g$ & Dividendo del próximo año, precio y tasa de crecimiento: $r_E=d_1/p+g$ (GORDON y SHAPIRO). & Bloque 2, anexo\\
$P_{i,0}$, $D_{i,1}$, $P_{i,1}$, $r_i$ & Precio de hoy, dividendo y precio al final del periodo, rentabilidad del título $i$. & Anexo\\
$w_i$ & Peso del título $i$ en una cartera. & Anexo\\
$\beta_{ik}$, $\lambda_k$ & Sensibilidad al factor $k$ y su prima (APT). & Anexo\\
$b_i$, $s_i$, $h_i$ & Cargas del modelo de tres factores. & Anexo\\
$B_1/P_0$ & Beneficio por acción del próximo año sobre el precio, inversa del PER. & Anexo\\
\end{longtable}

\section{3.B.4: Análisis del crecimiento de la empresa}
\begin{longtable}{SMD}
\cabecera
$V_B$ & Valor de la empresa $B$ por descuento de flujos. & Bloques 1-2\\
$FCL_t$, $BAII_t$, $\tau$ & Flujo de caja libre, beneficio antes de intereses e impuestos y tipo impositivo. & Bloques 1-2\\
$A_t$, $I_t$, $\Delta NOF_t$ & Amortizaciones, inversión y aumento de las necesidades operativas de fondos. & Bloques 1-2\\
$VT_T$, $k_B$ & Valor terminal y coste del capital del negocio de $B$. & Bloques 1-2\\
$E$ & Valor de mercado de la empresa, frente a la suma de sus partes (descuento del conglomerado). & Bloque 2\\
$S$, $s$ & Sinergia y sinergia percibida (SHLEIFER y VISHNY). & Bloque 2\\
$P$, $p$, $Q_A$, $Q_B$ & Precio por unidad de capital de $B$, fuerza negociadora de $B$ y valoraciones de $A$ y $B$: $P=pQ_A+(1-p)Q_B$. & Bloque 2\\
$N_A$ & Número de acciones de $A$ antes de la operación. & Bloque 2\\
$R_{it}$, $R_{mt}$, $\alpha_i$, $\beta_i$, $\varepsilon_{it}$ & Modelo de mercado de los estudios de acontecimientos: rentabilidad de la acción y del índice. & Anexo\\
$V$, $V_0$, $V^+$, $V^-$ & Valor del subyacente de la opción real: hoy y en los estados favorable y desfavorable. & Anexo\\
$u$, $d$, $q$ & Factores de subida y de bajada y probabilidad neutral al riesgo. & Anexo\\
$P$, $\delta$ & Precio de ejercicio y tasa de los flujos que reparte el subyacente. & Anexo\\
$C_0$, $C_u$, $C_d$ & Valor de la opción hoy y en cada estado. & Anexo\\
\end{longtable}

\section{3.A.10: Teoría de la demanda del consumidor (III)}
\begin{longtable}{SMD}
\cabecera
$x$, $\tilde x$ & Resultado de una lotería; resultado aleatorio. & Bloque 1, anexos\\
$p$, $p_n$ & Probabilidades de una lotería. & Bloque 1, anexos\\
$L$, $L'$, $L''$, $\bar L$, $\underline L$ & Loterías; la mejor y la peor. & Anexo de la utilidad esperada\\
$\alpha$ & Peso de la mezcla en el axioma de independencia. & Anexo de la utilidad esperada\\
$t$ & Peso de la mezcla en la demostración del teorema. & Anexo de la utilidad esperada\\
$\alpha_L$, $U(L)$ & Utilidad de la lotería $L$ (probabilidad de la mejor lotería indiferente) y utilidad esperada. & Anexo de la utilidad esperada\\
$\kappa$, $\eta$ & Coeficientes de una transformación afín de la utilidad, $\kappa U+\eta$ con $\kappa>0$. & Anexo de la utilidad esperada\\
$E[u]$ & Utilidad esperada. & Todo el tema\\
$EC$, $c(F,u)$ & Equivalente cierto. & Bloque 1, anexos\\
$\pi$, $\pi_p(x,\varepsilon,u)$ & Prima de riesgo (PRATT) y prima de probabilidad. & Bloque 1, anexos\\
$A(x)$, $R(x)$, $P(x)$ & Aversión absoluta y relativa al riesgo y prudencia absoluta, $-u'''/u''$. & Bloque 1, anexos\\
$\theta$ & Aversión relativa constante (CRRA): $u(x)=x^{1-\theta}/(1-\theta)$. & Bloque 1, anexos\\
$a$, $b$ & Coeficientes de las utilidades CARA, $-e^{-ax}/a$, y cuadrática, $x-\tfrac b2x^2$. & Anexo de la aversión\\
$w$ & Riqueza. & Bloque 1, anexos\\
$L$, $pL$, $\Pi$, $q$ & Pérdida, prima justa, prima de reserva, $\Pi=pL+\pi$, y precio del seguro por euro de cobertura. & Seguro\\
$\omega$, $\mu(\omega)$ & Estado de la naturaleza y probabilidad que le asigna el agente. & Información\\
$a$, $u(a,\omega)$, $w(a,\omega)$ & Acto y su utilidad y riqueza en cada estado. & Información\\
$s$, $q(s)$, $\mu(\omega\mid s)$ & Señal, su probabilidad y probabilidad revisada. & Información\\
$VIP$, $V$, $c^*$ & Valor de la información privada, valor de la información y precio máximo de la señal. & Información\\
$c_0$, $c_1$, $\tilde y$, $R_f$, $\delta$ & Consumo, renta futura incierta, tipo sin riesgo bruto y factor de descuento (ahorro precautorio). & Anexo de la aversión\\
$\sigma$ & Elasticidad de sustitución intertemporal, $1/\theta$ con la CRRA. & Anexo de la aversión\\
$v(x)$, $\alpha$, $\beta$, $\lambda$ & Función de valor de TVERSKY y KAHNEMAN, curvaturas de ganancias y pérdidas y aversión a las pérdidas. & Teoría de las perspectivas\\
$w(p)$, $\gamma$ & Función de ponderación de probabilidades y su parámetro. & Teoría de las perspectivas\\
$\pi_i^+$, $\pi_i^-$ & Pesos de decisión de las ganancias y de las pérdidas. & Teoría de las perspectivas\\
$C$ & Conjunto de probabilidades que el agente considera posibles (maxmin). & Ambigüedad\\
$\nu$, $E$, $f$, $U(f)$ & Capacidad, suceso, acto y utilidad de CHOQUET. & Ambigüedad\\
\end{longtable}

\section{3.B.25: Análisis de los instrumentos y de los mercados de derivados}
\begin{longtable}{SMD}
\cabecera
$S_t$, $S_T$ & Valor del subyacente en $t$ y al vencimiento, $T$. & Todo el tema\\
$f(S_T)$; $f$, $f_t$ & Pago al vencimiento, por unidad de subyacente o de nominal: la función que identifica el contrato; valor del derivado. & Todo el tema\\
$A$, $B$ & Parte larga (comprará el subyacente) y parte corta (lo venderá) del ejemplo conductor. & Todo el tema\\
$F$, $F_0$, $F_t$ & Precio a plazo pactado y precio a plazo en $0$ y en $t$. & Todo el tema\\
$K$ & Precio de ejercicio de una opción. & Todo el tema\\
$x^+$ & $\max\{x,0\}$. & Todo el tema\\
$L_{i-1}$, $L^F_i$, $\tau$, $s$, $s^*$, $m$ & Tipo variable de referencia fijado al principio del periodo $i$ y su tipo a plazo; duración de un periodo, en años; tipo fijo pactado en un swap o un FRA y tipo swap; diferencial de un préstamo sobre el tipo de referencia. & Todo el tema\\
$s$, $R$, $\gamma$ & Prima de un CDS (su pata fija), fracción que se recupera si hay impago e intensidad de impago. & Bloque 1 ($s$, $R$), anexos I y II\\
$b_t$ & Base, $S_t-F_t$. & Anexos I y II\\
$M_0$, $M_m$ & Garantía inicial y de mantenimiento de un futuro. & Anexo I\\
$w_1^A$, $w_2^A$, $w_1^B$, $w_2^B$, $u_A$, $u_B$, $\pi_1$, $\pi_2$, $e$, $e^*$ & Riqueza en los estados 1 y 2, utilidades y probabilidades del reparto del riesgo (caja de EDGEWORTH, condición de BORCH); dotación y asignación tras el contrato. & Bloque 1, anexo I\\
$V(w)$, $Y(I)$, $C(\cdot)$, $I$, $w$ & Valor de la empresa, rendimiento de la inversión, coste de la financiación externa, inversión y caja interna (FROOT, SCHARFSTEIN y STEIN, 1993). & Anexo I\\
$r_f$, $r_f^*$, $r_a$, $r_T$ & Tipo sin riesgo con capitalización continua, nacional y extranjero; tipo con capitalización anual, $(1+r_a)^t=e^{r_ft}$, y tipo del periodo hasta $T$. & Anexo II\\
$B(0,t)$, $B_{\text{fijo}}$, $B_{\text{var}}$, $B_i$ & Precio hoy de un bono sin riesgo que paga $1$ en $t$; valor de los bonos fijo y variable de un swap; precio del bono entregable $i$ de un futuro. & Anexos I y II\\
$\delta$, $a$, $y$, $I$ & Rendimiento continuo del subyacente, coste continuo de almacenamiento, rendimiento de conveniencia y valor actual de las rentas. & Anexo II\\
$c$, $p$, $C$, $P$ & Primas o precios de la \textit{call} y de la \textit{put} europeas y americanas. & Todo el tema\\
$\mu$, $\sigma$, $W_t$, $\varepsilon$ & Rentabilidad esperada y volatilidad del subyacente, movimiento browniano y variable normal estándar. & Bloque 2 ($\mu$, $\sigma$), anexo II\\
$u$, $d$, $q$, $\tau$, $Z$ & Factores de subida y de bajada del árbol, probabilidad neutral al riesgo, duración del periodo y cantidad en el activo sin riesgo de la réplica. & Bloque 2 ($q$, $\tau$, $Z$), anexo II\\
$S_u$, $S_d$; $S_{uu}$, $S_{ud}$, $S_{dd}$; $f_0$, $f_u$, $f_d$; $\pi$ & Precio del subyacente al vencimiento si sube o si baja, en el mundo de dos estados, y, en el árbol de dos periodos, tras dos subidas, una de cada o dos bajadas; precio del derivado hoy y su pago en cada estado; probabilidad real de subida. & Bloque 2\\
$d_1$, $d_2$, $N(\cdot)$, $N'(\cdot)$ & Argumentos de BLACK-SCHOLES y función de distribución y de densidad normales. & Todo el tema\\
$\Delta$, $\Gamma$, $\Theta$, $\nu$, $\rho$; $\Delta_0$, $\Delta_u$, $\Delta_d$; $\Delta_j$ & Griegas: sensibilidades al subyacente, curvatura, paso del tiempo, volatilidad y tipo de interés; delta de la réplica en cada nudo del árbol; pendiente de un pago lineal a tramos. Las variaciones se escriben $dS$, $dt$ y no con $\Delta$. & Bloque 2, anexos I y II\\
$\mathbb{Q}$, $\mathbb{P}$ & Medida neutral al riesgo y medida real. & Todo el tema\\
$\psi_\omega$, $q_\omega$, $\pi_\omega$; $x_{\omega j}$, $V_j$, $h_j$ & Precio de estado y probabilidades neutral al riesgo y real del estado $\omega$; pago y precio del activo $j$ y unidades en cartera. El factor de descuento estocástico es $\psi_\omega/\pi_\omega$. & Bloque 2 ($\psi_\omega$, $q_\omega$), anexo II\\
$\lambda$ & Precio de mercado del riesgo, $(\mu-r_f)/\sigma$. & Anexo II\\
$\eta$, $\bar J$, $\sigma_J$ & Intensidad de los saltos, variación media del precio en cada salto y su desviación típica (MERTON, 1976). & Anexo II\\
$h^*$, $h$, $\theta$ & Razón de cobertura de mínima varianza, unidades de los instrumentos de cobertura y aversión al riesgo. & Anexo II\\
$V$, $V_t$; $V_0$, $D$ & Valor de un contrato, de un activo o de una cartera; en el modelo de MERTON (1974), valor de los activos y nominal de la deuda, con las acciones como una \textit{call}, $c$. & Anexo II\\
$x$ & Base de divisas: desviación de la paridad cubierta; en el cuerpo, una cantidad genérica (un movimiento del precio, la pérdida de una cobertura). & Anexo II; cuerpo\\
\end{longtable}

\section{3.B.24: Análisis de los instrumentos financieros de renta fija}
\begin{longtable}{SMD}
\cabecera
$P$, $N$, $c$, $T$ & Precio, nominal, tasa del cupón (el cupón es $cN$; $C$ queda para la convexidad) y vencimiento de un bono. & Todo el tema\\
$CF_t$, $m$, $a_{T,r}$ & Pago de la fecha $t$, número de cupones al año y factor de anualidad, el valor actual de una renta de una unidad durante $T$ periodos al tipo $r$. & Todo el tema; $m$, anexo I\\
$r$, $y$ & Tipo único con que se descuentan todos los pagos (bloque 1) y rendimiento al vencimiento o TIR. & Todo el tema\\
$y_c$, $y_p$, $i_d$, $d$ & Rendimiento corriente, $cN/P$; rendimiento por periodo de cupón; tipo de descuento de las letras de Estados Unidos y días hasta su vencimiento. & Bloque 1, anexo I (I.3)\\
$P^{\text{sucio}}$, $P^{\text{limpio}}$, $\tau$, $k$, $n$ & Precio con y sin el cupón corrido, fracción transcurrida del periodo de cupón, índice de los pagos pendientes y número de pagos (Orden ECM/2/2026, art. 14.2). & Anexo I (I.2)\\
$y^{\text{real}}$, $\pi$, $\pi^e$ & Rendimiento real, inflación e inflación esperada (FISHER). & Anexo I (I.3)\\
$H$, $V_H$ & Horizonte del inversor y valor en $H$ de la inversión con los pagos reinvertidos. & Bloque 1, anexo I (I.3, I.8)\\
$r_f$, $s$, $p$, $\ell$, $\lambda$ & Tipo sin riesgo; diferencial, $s=y-r_f\approx p\,\ell$; probabilidad de impago; pérdida en caso de impago, en fracción de lo prometido (en el modelo de intensidad, del valor de mercado); intensidad de impago. & Bloque 1, anexo I (I.4)\\
$h$, $y^*$, $H_1<H_2<H_3$ & Horquilla; rendimiento de un bono perfectamente líquido; horizontes de las clientelas de AMIHUD y MENDELSON. & Bloque 2, anexo I (I.5)\\
$y_{\text{exento}}$, $\tau$ & Rendimiento de un bono exento y tipo impositivo. & Anexo I (I.5)\\
$P_v$, $P_c$, $V$ & Precios a los que vende y compra el intermediario y valor del bono. & Bloque 2\\
$D$, $w_t$, $D^*$ & Duración de MACAULAY y peso de cada pago en el precio; duración modificada, $D/(1+y)$. & Bloque 1, anexo I (I.6-I.8)\\
$C$, $M^2$, $R$ & Convexidad, dispersión de los pagos alrededor de la duración y rentabilidad de un periodo, $R\approx y\,\Delta t-D^*\Delta y+\tfrac12 C(\Delta y)^2$. & Bloque 1, anexo I (I.7)\\
$A$, $L$, $E$, $D_A^*$, $D_L^*$ & Activos, pasivos y capital de un balance, $E=A-L$, y duraciones modificadas de los activos y de los pasivos. & Bloque 1, anexo I (I.8)\\
$B(0,t)$, $r_t$, $r_n$ & Precio hoy de un cupón cero sin riesgo que paga una unidad en $t$ y tipo al contado a $t$ o a $n$ años. & Bloque 3, anexo I (I.2, I.9-I.14)\\
$f_{t,s}$, $c_n$ & Tipo a plazo entre $t$ y $s$ y tipo par a $n$ años. & Bloque 3, anexo I (I.9)\\
$n$, $\beta_0$, $\beta_1$, $\beta_2$, $\tau$ & Plazo, en años; nivel, pendiente y curvatura de NELSON y SIEGEL, y plazo en que actúa la curvatura. & Bloque 3, anexo I (I.10)\\
$i_k$, $E[i_k]$, $f_{k,k+1}$, $\phi_n$ & Tipo a un año dentro de $k$ años, su esperanza, el tipo a plazo de ese año y prima por plazo a $n$ años. & Bloque 3, anexo I (I.12-I.13)\\
$S_0$, $S_1$, $\phi_0$, $\phi_1$ & Oferta neta de deuda a largo plazo y prima por plazo antes y después de una emisión (VAYANOS y VILA). & Bloque 3\\
$m_{t+1}$, $B_t^{(n)}$ & Factor de descuento estocástico y precio en $t$ del cupón cero a $n$ años. & Anexo I (I.14)\\
$i$, $\bar\imath$, $\kappa$, $\sigma$, $W$ & Tipo corto, su nivel normal, velocidad con que vuelve a él, volatilidad y movimiento browniano (VASICEK). & Bloque 3, anexo I (I.14)\\
$a_n$, $b_n$, $r_\infty$ & Coeficientes del tipo a $n$ años en el modelo de VASICEK, $r_n=a_n+b_n\,i$, y tipo a muy largo plazo. & Bloque 3, anexo I (I.14)\\
\end{longtable}

\section{3.B.13: Mercado de divisas. Operaciones e instrumentos}
\begin{longtable}{SMD}
\cabecera
$A$, $B$, $D$, $E$ & Agentes del ejemplo conductor: la empresa que cobrará dólares, la que tendrá que pagarlos, el intermediario y el banco o la empresa de la zona del euro que se financia en dólares. & Todo el tema\\
$S$, $S_0$, $S_T$, $F$ & Tipo de cambio al contado en euros por dólar (cotización directa: sube cuando el euro se deprecia), hoy y al vencimiento, y tipo a plazo con el mismo convenio. & Todo el tema\\
$T$, $\tau$, $N$ & Vencimiento; plazo en fracción de año (días/360 en euros y en dólares, /365 en libras), y nominal en dólares. & Todo el tema\\
$r_f$, $r_f^*$ & Tipos sin riesgo del euro y del dólar, con el interés simple del mercado monetario. & Todo el tema\\
$\mathrm{EUR/USD}$ & Cotización de mercado del par, en dólares por euro: $S=1/\mathrm{EUR/USD}$. & Bloque 1, anexo I\\
$b$, $a$, $m$, $h$ & Precios comprador y vendedor del par ($b_{\mathrm{EUR/USD}}$, $a_{\mathrm{EUR/USD}}$), precio medio y horquilla relativa. & Anexos I (I.2) y II (II.1)\\
$\mathrm{TCEN}_t$, $c_{j,t}$, $w_j$ & Tipo de cambio efectivo nominal, unidades de la moneda $j$ por unidad de la nacional y peso del socio $j$. & Anexo I (I.2)\\
$b^{d}$, $a^{d}$ & Precios directos del par en el arbitraje triangular. & Anexo I (I.9)\\
$S^{b}$, $S^{a}$, $F^{b}$, $F^{a}$; $r^{p}$, $r^{d}$, $r^{p*}$, $r^{d*}$ & Contado y plazo, comprador y vendedor; tipos a los que se pide prestado y se deposita, en euros y en dólares: la banda de la paridad cubierta. & Anexo I (I.9)\\
$F^{L}$, $S^{L}_T$, $r^{L}$ & Tipo de un contrato a plazo sin entrega y fijación al vencimiento, en moneda local por dólar, y tipo de interés implícito de la moneda local. & Anexo I (I.4)\\
$r^{\$}_{\mathrm{sint}}$, $x$ & Tipo sintético de los dólares obtenidos con un \textit{FX swap}, $\frac{1}{\tau}\left[\frac{S}{F}(1+r_f\tau)-1\right]$, y base de divisas, $x=r_f^*-r^{\$}_{\mathrm{sint}}$. & Bloques 1 y 3, anexos I y II\\
$c$, $c^*$ & Cupones en euros y en dólares de un swap de divisas. & Anexo I (I.5, figura)\\
$c$, $p$, $K$, $\Phi$, $\sigma$ & Primas de la \textit{call} y de la \textit{put}, precio de ejercicio, distribución normal estándar (no $N$, que es el nominal) y volatilidad, en GARMAN y KOHLHAGEN, que usa la capitalización continua. & Anexo I (I.6)\\
$\sigma_{ATM}$, $\sigma_{25c}$, $\sigma_{25p}$, $RR_{25}$, $BF_{25}$ & Volatilidades cotizadas en el dinero y de las opciones de delta 25, \textit{risk reversal} y mariposa. & Anexo I (I.6)\\
$\beta_S$, $V_t$ & Exposición cambiaria de ADLER y DUMAS (no $\beta$, que es la beta de finanzas) y valor de la empresa. & Anexo I (I.7)\\
$z_{t+1}$; $r_{\mathrm{alto}}$, $r_{\mathrm{bajo}}$ & Rentabilidad del \textit{carry trade}, $z_{t+1}\approx(r_{f,t}^*-r_{f,t})+\Delta\ln S_{t+1}$; tipos de interés de la moneda en que se invierte y de la moneda en que se endeuda. & Anexo I (I.8); bloque 2\\
$q$, $m$, $\theta$, $\sigma^2$ & Inventario del intermediario, en unidades de la divisa base; valor que espera de la divisa; aversión al riesgo, y varianza del tipo de cambio mientras mantiene la posición. & Anexo II (II.2)\\
$V$, $V_H$, $V_L$, $\mu$ & Valor de la divisa, alto y bajo, y fracción de clientes informados (GLOSTEN y MILGROM). & Anexo II (II.3)\\
$v$, $p_0$, $\sigma_v$, $u$, $\sigma_u$, $x$, $y$, $p$, $\lambda$, $\kappa$ & Valor de la divisa, precio inicial y su incertidumbre; flujo de los clientes no informados y la suya; orden del informado, $x=\kappa(v-p_0)$; flujo total, $y=x+u$; precio, $p=p_0+\lambda y$; impacto de una unidad de flujo, $\lambda=\sigma_v/(2\sigma_u)$, e intensidad con que negocia el informado, $\kappa=\sigma_u/\sigma_v$ (KYLE, 1985). & Anexo II (II.3)\\
$y_t$, $\gamma_1$, $\gamma_2$, $\varepsilon_t$ & Flujo de órdenes y coeficientes de EVANS y LYONS, $\Delta\ln S_t=\gamma_1\Delta(r_{f,t}-r_{f,t}^*)+\gamma_2y_t+\varepsilon_t$. & Anexo II (II.5)\\
$s_t$, $f_t$, $\alpha$, $\beta$, $\pi_t$, $\mathbb{E}_t$ & Logaritmos del contado y del plazo; constante y coeficiente de la regresión de FAMA; prima de riesgo de cambio, $\pi_t\equiv f_t-\mathbb{E}_t[s_{t+1}]$, y esperanza con la información de $t$. & Anexo II (II.7)\\
$h_{12}$, $h_{13}$, $h_{23}$; $M_1$, $M_2$, $M_3$; $n$ & Horquillas de los mercados bilaterales entre la moneda vehicular, $M_1$, y las monedas $M_2$ y $M_3$; número de monedas, con $n(n-1)/2$ mercados posibles. & Bloque 3, anexo III (III.7)\\
\end{longtable}

\end{document}
